% Part: computability
% Chapter: computability-theory
% Section: equiv-ce-defs

\documentclass[../../../include/open-logic-section]{subfiles}

\begin{document}

\olfileid[pt]{cmp}{thy}{eqc}

\olsection[Definições de conjuntos c.e.]{Definições equivalentes de
  conjuntos computavelmente enumeráveis}


O resultado a seguir apresenta vários enunciados equivalentes
importantes do que significa ser computavelmente enumerável.

\begin{thm}
\ollabel{thm:ce-equiv}
Seja $S$ um conjunto de números naturais. Então as seguintes afirmações
são equivalentes:
\begin{enumerate}
\item\ollabel{case:ce} $S$ é computavelmente enumerável.
\item\ollabel{case:ran-pc} $S$ é a imagem de uma função computável \emph{parcial}.
\item\ollabel{case:ran-prim} $S$ é vazio ou é a imagem de uma função recursiva primitiva.
\item\ollabel{case:ce-domain} $S$ é o \emph{domínio} de uma função parcialmente computável.
\end{enumerate}
\end{thm}

\begin{explain}
As três primeiras cláusulas dizem que podemos, de modo equivalente,
tomar qualquer conjunto computavelmente enumerável não vazio como
enumerado por uma função computável, por uma função parcialmente
computável ou por uma função recursiva primitiva. A quarta cláusula nos
diz que, se $S$ é computavelmente enumerável, então, para algum
índice~$e$,
\[
S = \Setabs{x}{\cfind{e}(x) \fdefined}.
\]
Em outras palavras, $S$ é o conjunto das entradas para as quais a
computação de $\cfind{e}$ para. Por essa razão, conjuntos
computavelmente enumeráveis são às vezes chamados de
\emph{semidecidíveis}: se um número está no conjunto, você acaba
recebendo um ``sim'', mas, se ele não está, nunca recebe um ``não''!{}
\end{explain}

\begin{proof}
Como toda função recursiva primitiva é computável e toda função
computável é parcialmente computável, \olref{case:ran-prim}
implica \olref{case:ce}, e \olref{case:ce}
implica~\olref{case:ran-pc}. (Observe que, se $S$ for vazio, $S$~será
a imagem da função parcialmente computável que não está definida em
ponto algum.) Se mostrarmos que \olref{case:ran-pc} implica
\olref{case:ran-prim}, teremos mostrado que as três primeiras cláusulas
são equivalentes.

Suponha, então, que $S$ seja a imagem da função parcialmente computável
$\cfind{e}$. Se $S$ for vazio, terminamos. Caso contrário, seja $a$ um
elemento qualquer de~$S$. Pelo teorema da forma normal de Kleene,
podemos escrever
\[
\cfind{e}(x) = U(\umin{s}{T(e, x, s)}).
\]
Em particular, $\cfind{e}(x) \fdefined$ e $= y$ se, e somente se,
existe um $s$ tal que $T(e, x, s)$ e $U(s) = y$. Defina $f(z)$ por
\[
f(z) = \begin{cases}
  U((z)_1) & \text{se $T(e, (z)_0, (z)_1)$} \\
  a        & \text{caso contrário.}
\end{cases}
\]
Então $f$ é recursiva primitiva, pois $T$ e $U$
o são. \iftag{TMs}{Em termos de máquinas de Turing, se $z$ codifica
  um par $\tuple{(z)_0, (z)_1}$ tal que $(z)_1$ é uma computação que
  para da máquina~$M_e$ na entrada $(z)_0$, então $f$ devolve a saída
  da computação; caso contrário, devolve~$a$.}

Precisamos mostrar que $S$ é a imagem de~$f$, O.s.\@, que, para todo
número natural~$y$, $y \in S$ se, e somente se, esse número pertence
à imagem de~$f$. No sentido direto, suponha que $y \in S$. Então $y$ pertence
à imagem de $\cfind{e}$, de modo que, para algum $x$ e algum~$s$,
$T(e,x,s)$ vale e $U(s) = y$. Mas então
$y = f(\tuple{x,s})$. Reciprocamente, suponha que $y$ pertença à
imagem de~$f$. Então ou $y = a$, ou, para algum~$z$,
$T(e,(z)_0,(z)_1)$ e $U((z)_1) = y$. Como, neste último caso,
$\cfind{e}(x) \fdefined = y$, em qualquer dos casos $y$ pertence
a~$S$.

(A notação $\cfind{e}(x) \fdefined = y$ significa
``$\cfind{e}(x)$ está definida e é igual a $y$.'' Poderíamos
simplesmente usar $\cfind{e}(x) = y$, mas a seta adicional às vezes
ajuda a nos lembrar de que estamos tratando de uma função parcial.)

Para concluir a demonstração de \olref{thm:ce-equiv}, basta mostrar que
\olref{case:ce} e~\olref{case:ce-domain} são equivalentes. Primeiro,
mostremos que \olref{case:ce} implica~\olref{case:ce-domain}.
Suponha que $S$ seja a imagem de uma função computável~$f$, O.s.\@,
\[
S = \Setabs{y}{\text{para algum $x$, } f(x) = y}.
\]
Seja
\[
g(y) = \umin{x}{(f(x) = y)}.
\]
Então $g$ é uma função parcialmente computável, e $g(y)$ está definida
se, e somente se, para algum~$x$, $f(x) = y$. Em outras palavras, o
domínio de $g$~é a imagem de~$f$. \iftag{TMs}{Em termos de máquinas de
Turing: dada uma máquina de Turing~$F$ que enumera os elementos de~$S$,
seja $G$ a máquina de Turing que semidecide $S$ procurando entre as
saídas de~$F$ para verificar se um dado elemento pertence ao conjunto,
para se ele pertencer e continua procurando para sempre se ele não
pertencer.}{}

Por fim, para mostrar que \olref{case:ce-domain} implica~\olref{case:ce},
suponha que $S$~seja o domínio da função parcialmente computável
$\cfind{e}$, O.s.\@,
\[
S = \Setabs{x}{\cfind{e}(x) \fdefined}.
\]
Se $S$ for vazio, terminamos; caso contrário, seja $a$ um elemento
qualquer de~$S$. Defina $f$ por
\[
f(z) = \begin{cases}
(z)_0 & \text{se $T(e,(z)_0,(z)_1)$} \\
a & \text{caso contrário.}
\end{cases}
\]
Então, como acima, um número $x$ pertence à imagem de~$f$ se, e
somente se, $\cfind{e}(x) \fdefined$, O.s.\@, se, e somente se,
$x \in S$. \iftag{TMs}{Em termos de máquinas de Turing: dada uma
máquina $M_e$ que semidecide $S$, enumere os elementos de $S$
percorrendo todas as computações possíveis de máquinas de Turing e
devolvendo as entradas que correspondem a computações que param.}{}
\end{proof}

A cláusula~\olref{case:ce-domain} de \olref{thm:ce-equiv} nos fornece
uma maneira conveniente de enumerar os conjuntos computavelmente
enumeráveis: para cada~$e$, seja $W_e$ o domínio de $\cfind{e}$, O.s.\@,
\[
W_e = \Setabs{x}{\cfind{e}(x) \fdefined}.
\] 
Então, se $A$ é um conjunto computavelmente enumerável qualquer,
$A = W_e$ para algum~$e$.

O resultado a seguir fornece ainda outra caracterização dos conjuntos
computavelmente enumeráveis.

\begin{thm}
\ollabel{thm:exists-char}
Um conjunto $S$ é computavelmente enumerável se, e somente se, existe
uma relação computável $R(x,y)$ tal que
\[
S = \Setabs{ x }{ \lexists[y][R(x,y)] }.
\]
\end{thm}

\begin{proof}
No sentido direto, suponha que $S$ seja computavelmente enumerável.
Então, para algum $e$, $S = W_e$. Para esse valor de~$e$, podemos
escrever $S$ como
\[
S = \Setabs{ x }{ \lexists[y][T(e, x, y)] }.
\]
No sentido inverso, suponha que $S = \Setabs{ x }{
  \lexists[y][R(x, y)] }$. Defina $f$~por
\[
f(x) \simeq \umin{y}{R(x, y)}.
\]
Então $f$ é parcialmente computável, e $S$ é o domínio de~$f$.
\end{proof}


\end{document}
